% ======================================================================
%  EduFocus | Méthodes statistiques, ML et data mining — un cours appliqué
%  Compilation : pdflatex (3 passes pour la table des matières)
% ======================================================================
\documentclass[11pt,a4paper]{article}

\usepackage[utf8]{inputenc}
\usepackage[T1]{fontenc}
\usepackage[french]{babel}
\usepackage[a4paper,margin=2.4cm]{geometry}
\usepackage{amsmath,amssymb,amsthm}
\usepackage{booktabs}
\usepackage{xcolor}
\usepackage{tcolorbox}
\tcbuselibrary{breakable,skins}
\usepackage{listings}
\usepackage{graphicx}
\usepackage[hidelinks]{hyperref}
\usepackage{tabularx}
\usepackage{enumitem}

% ---------- Palette EduFocus ----------
\definecolor{eduTeal}{HTML}{4EC3A3}
\definecolor{eduAmber}{HTML}{EEB74F}
\definecolor{eduRed}{HTML}{EF6F5F}
\definecolor{eduDark}{HTML}{26323A}
\definecolor{eduMute}{HTML}{A19077}

% ---------- Boîtes ----------
\newtcolorbox{meth}[1]{
  colback=white, colframe=eduDark, boxrule=0.9pt,
  title={\textbf{#1}}, fonttitle=\color{white}, coltitle=white,
  colbacktitle=eduDark, breakable, sharp corners=south
}
\newtcolorbox{where}{
  colback=eduTeal!6, colframe=eduTeal, boxrule=0.7pt,
  title={\textbf{Où c'est utilisé dans EduFocus}},
  fonttitle=\bfseries, coltitle=black!75, colbacktitle=eduTeal!40,
  breakable, sharp corners=south
}
\newtcolorbox{notebox}{
  colback=eduAmber!10, colframe=eduAmber, boxrule=0.7pt,
  breakable, sharp corners=south
}

\newcommand{\txt}[1]{\texttt{#1}}

% ---------- Listings ----------
\lstset{
  language=Python, basicstyle=\ttfamily\small,
  keywordstyle=\color{eduDark}\bfseries,
  commentstyle=\color{eduMute}\itshape,
  stringstyle=\color{eduRed},
  frame=single, rulecolor=\color{black!25}, backgroundcolor=\color{black!3},
  breaklines=true, columns=fullflexible,
  morekeywords={as,None,True,False,def,return}
}

% ---------- En-têtes ----------
\title{\vspace{-1.2cm}
  {\color{eduDark}\textbf{EduFocus}}\\[0.15cm]
  {\Large Méthodes statistiques, apprentissage machine et data mining}\\[0.25cm]
  {\large Un cours appliqué : les mathématiques derrière l'analyse\\ de l'exclusion scolaire en Mauritanie}
}
\author{Projet EduFocus -- Datathon IndabaX Mauritanie 2026}
\date{}

\begin{document}

\maketitle
\vspace{-1.1cm}
\begin{center}
\begin{minipage}{0.92\textwidth}
\small\itshape\color{eduMute}
Ce document est un \emph{cours} qui explique \textbf{pourquoi} et \textbf{comment}
chaque méthode est utilisée dans EduFocus : d'abord l'idée, ensuite la
\textbf{formule mathématique}, puis \textbf{où} elle intervient dans le code,
l'API et l'interface.
\end{minipage}
\end{center}

\tableofcontents
\newpage

% ======================================================================
\section{Introduction : le problème et les trois familles de méthodes}
% ======================================================================

EduFocus analyse l'exclusion scolaire en Mauritanie : \textbf{375\,566 enfants
de 6--14 ans} (soit \textbf{33,1\,\%}) sont hors de l'école formelle selon
l'enquête EPCV 2019, sur une population cible de 1\,135\,657 enfants répartis
dans les \textbf{13 wilayas}.

Pour transformer ces données brutes en décisions d'investissement, le projet
utilise trois familles de méthodes complémentaires :

\begin{enumerate}[leftmargin=2.2em]
  \item \textbf{Statistiques} : décrire (moyennes, médianes, concentration) et
        mesurer les liaisons entre indicateurs (corrélation de Pearson).
  \item \textbf{Apprentissage machine} : \emph{non supervisé} pour regrouper
        les wilayas (K-means, silhouette) et \emph{supervisé} pour expliquer
        les déterminants individuels (régression logistique).
  \item \textbf{Data mining et graphes} : règles d'association (Apriori) et
        analyse de réseaux (détection de communautés de Louvain) pour
        découvrir des structures non évidentes.
\end{enumerate}

\begin{notebox}
\textbf{Fil rouge.} Chaque méthode est décrite dans ce cours comme suit :
\emph{idée} $\to$ \emph{formule} $\to$ \emph{exemple chiffré} $\to$
\emph{où c'est utilisé}. Un tableau récapitulatif complet est donné au
§~\ref{sec:resume}.
\end{notebox}

% ======================================================================
\section{Statistiques descriptives et inférentielles}
% ======================================================================

% ----------------------------------------------------------------------
\subsection{Normalisation min--max}
\label{sec:minmax}

\textbf{Idée.} Les indicateurs n'ont pas la même unité (un taux en \%, un
effectif en milliers, un ratio sans dimension). Pour les comparer et les
combiner, on les ramène sur une échelle commune $[0,\,100]$.

\begin{meth}{Formule}
\[
x'_i \;=\; \frac{x_i - \min(x)}{\max(x) - \min(x)} \times 100
\]
où $\min(x)$ et $\max(x)$ sont les valeurs extrêmes observées sur les
13 wilayas. Si $\max(x)=\min(x)$ (pas de variation), on fixe $x'_i=50$.
\end{meth}

\begin{where}
\txt{backend/analytics/index.py} -- fonction \txt{normalize()} : normalise
\textbf{le volume} (nombre d'enfants hors école, en échelle log, cf.~\ref{sec:log})
et \textbf{l'intensité} (taux de hors-école) avant de construire l'indice IPE
(§~\ref{sec:ipe}). Accessible via \txt{/api/wilayas} et affiché sur la page
\txt{/wilayas}, la carte \txt{/carte} et le rapport \txt{/rapport}.
\end{where}

% ----------------------------------------------------------------------
\subsection{Standardisation z-score}
\label{sec:zscore}

\textbf{Idée.} Pour un \emph{clustering}, les données doivent être centrées et
réduites : chaque variable doit avoir la même échelle, sinon les variables
grandes en valeur domineraient les calculs de distance.

\begin{meth}{Formule}
\[
z_i \;=\; \frac{x_i - \mu}{\sigma}, \qquad
\mu=\frac{1}{n}\sum_{i=1}^{n}x_i, \qquad
\sigma=\sqrt{\frac{1}{n}\sum_{i=1}^{n}(x_i-\mu)^2}.
\]
Après transformation : moyenne $0$, écart-type $1$.
\end{meth}

\begin{where}
\txt{backend/analytics/clustering.py} -- \txt{StandardScaler().fit\_transform(...)}
sur 8 features (intensité, pauvreté, ruralité, dépendance jeunes, mahadra,
aucune instruction, écart de genre, densité d'écoles) avant le K-means
(§~\ref{sec:kmeans}).
\end{where}

% ----------------------------------------------------------------------
\subsection{Transformation logarithmique}
\label{sec:log}

\textbf{Idée.} Le nombre d'enfants hors école varie de quelques milliers
(Nouadhibou) à des centaines de milliers (Hodh Ech Chargui). Une échelle
logarithmique « compresse » les grands nombres et rend la variable symétrique.

\begin{meth}{Formule}
\[
y \;=\; \log(1+x) \qquad\text{(log1p)}
\]
Le $+1$ garantit que $x=0 \mapsto y=0$. Utilisée avant la normalisation
min--max du \emph{volume} et dans la matrice de stratégie.
\end{meth}

\begin{where}
\txt{backend/analytics/index.py} (\txt{np.log1p}) et
\txt{backend/analytics/matrice.py} (axe « volume » du scatter). Page
\txt{/strategies}.
\end{where}

% ----------------------------------------------------------------------
\subsection{Médiane et partage en quadrants}
\label{sec:median}

\textbf{Idée.} La médiane coupe une distribution en deux moitiés égales. Elle
est robuste aux valeurs extrêmes (contrairement à la moyenne).

\begin{meth}{Définition}
Pour un échantillon trié $x_{(1)}\le\dots\le x_{(n)}$ :
\[
\text{médiane} \;=\;
\begin{cases}
x_{(\frac{n+1}{2})} & \text{si } n \text{ impair}\\[0.25em]
\frac{1}{2}\big(x_{(\frac{n}{2})}+x_{(\frac{n}{2}+1)}\big) & \text{si } n \text{ pair}.
\end{cases}
\]
\end{meth}

\begin{where}
\txt{backend/analytics/matrice.py} : chaque wilaya est classée \emph{haut/bas}
selon la médiane du volume log et la médiane de l'intensité $\to$ \textbf{4
quadrants} (« Effort massif », « Ciblage des effectifs », « Traitement
local », « Consolidation »). Page \txt{/strategies}, endpoint \txt{/api/matrice}.
\end{where}

% ----------------------------------------------------------------------
\subsection{Corrélation de Pearson}
\label{sec:pearson}

\textbf{Idée.} Mesure la force et le sens de la \emph{liaison linéaire} entre
deux indicateurs, en unités normalisées (donc comparable entre paires).

\begin{meth}{Formule}
\[
r \;=\;
\frac{\sum_{i=1}^{n}(x_i-\bar{x})(y_i-\bar{y})}
{\sqrt{\sum_{i=1}^{n}(x_i-\bar{x})^2}\;\sqrt{\sum_{i=1}^{n}(y_i-\bar{y})^2}}
\;=\;\frac{\operatorname{Cov}(x,y)}{\sigma_x\,\sigma_y}
\]
Propriétés : $r\in[-1,\,1]$ ; $r=1$ (resp.\ $-1$) : corrélation linéaire
positive (resp.\ négative) parfaite ; $r=0$ : aucune liaison linéaire.
\end{meth}

\begin{notebox}
\textbf{Ce que $r$ mesure et ne mesure pas.} La corrélation est \emph{symétrique}
et \emph{sans direction de causalité} : si le taux de pauvreté et le taux de
hors-école corrèlent à $r=0{,}84$, cela ne prouve pas que la pauvreté
\emph{cause} le hors-école. C'est le rôle de la régression logistique
(§~\ref{sec:logit}) d'approcher la causalité en \emph{ajustant} les facteurs.
\end{notebox}

\begin{where}
Trois usages :
\begin{itemize}[leftmargin=2em]
  \item \txt{backend/analytics/indicateurs.py} : matrice complète
        $12\times12$ affichée en heatmap sur \txt{/indicateurs} et
        \txt{/rapport} (endpoint \txt{/api/indicateurs}). Paires les plus
        fortes :
        dépendance jeunes $\times$ part 0--14 ans ($r=1{,}00$), aucune
        instruction $\times$ adultes 15+ sans instruction ($r=0{,}94$),
        ruralité $\times$ dépendance jeunes ($r=0{,}91$).
  \item \txt{backend/analytics/graph.py} : réseau de similarité entre
        wilayas (profils centrés-réduits, $|r|>0{,}70$, top-2 voisins) et
        graphe de corrélations entre indicateurs ($|r|>0{,}75$). Page
        \txt{/reseau}.
  \item \txt{backend/analytics/concentration.py} : n'utilise pas $r$ mais
        quantifie la concentration (voir §~\ref{sec:gini}).
\end{itemize}
\end{where}

% ----------------------------------------------------------------------
\subsection{Concentration : Pareto, Lorenz et indice de Gini}
\label{sec:gini}

\textbf{Idée.} Le hors-école n'est pas uniformément réparti. On veut savoir
\textbf{combien de wilayas portent la moitié du problème} et mesurer
l'inégalité de la répartition (loi de Pareto « 20\,\% des zones $\to$
80\,\% du problème »).

\begin{meth}{Formules}
\textbf{Courbe de Lorenz} (part cumulée) : on trie $x_{(1)}\ge\dots\ge x_{(n)}$
par ordre décroissant, alors
\[
\text{Lorenz}(k) \;=\; \frac{\sum_{i=1}^{k}x_{(i)}}{\sum_{i=1}^{n}x_{(i)}} \times 100.
\]
\textbf{Indice de Gini} (formule par paires, utilisée dans le code) :
\[
G \;=\; \frac{\sum_{i=1}^{n}\sum_{j=1}^{n}|x_i-x_j|}
{2\,n\sum_{i=1}^{n}x_i}.
\]
Interprétation : $G=0$ répartition parfaitement égale, $G=1$ concentration
totale. $G$ est lié à l'aire entre la courbe de Lorenz et la diagonale
(égalité parfaite).
\end{meth}

\begin{where}
\txt{backend/analytics/concentration.py} : calcule la part cumulée des 5
premières wilayas (\textbf{top5}), le nombre de wilayas pour atteindre
50\,\% du problème, la courbe de Lorenz et le Gini. Endpoints
\txt{/api/concentration} ; affiché sur \txt{/} (Home), \txt{/strategies}
et \txt{/rapport}.
\end{where}

% ----------------------------------------------------------------------
\subsection{Tendance et régression linéaire simple (moindres carrés)}
\label{sec:ols}

\textbf{Idée.} Pour projeter le taux de hors-école jusqu'en 2030, on mesure la
pente de la tendance passée (World Bank, 2008--2024) par régression linéaire.

\begin{meth}{Formule}
Modèle $y = \beta_0 + \beta_1 x + \varepsilon$. Le moindre carré minimise
$\sum_i\big(y_i - \beta_0 - \beta_1 x_i\big)^2$, solution :
\[
\beta_1 \;=\; \frac{\sum_{i}(x_i-\bar{x})(y_i-\bar{y})}
{\sum_{i}(x_i-\bar{x})^2}, \qquad
\beta_0 \;=\; \bar{y} - \beta_1\bar{x}.
\]
\end{meth}

\begin{where}
\txt{backend/analytics/scenarios.py} -- \txt{np.polyfit(x, y, 1)} : pente
annuelle bornée dans $[-1{,}0,\,0]$ points par an, prolongée sur
$\text{2030}-2022$ ans. Endpoint \txt{/api/scenarios}, page \txt{/strategies}.
Les séries brutes WDI sont servies par \txt{/api/trends} sur \txt{/tendances}.
\end{where}

% ======================================================================
\section{Apprentissage machine non supervisé : typologie des wilayas}
% ======================================================================

% ----------------------------------------------------------------------
\subsection{K-means}
\label{sec:kmeans}

\textbf{Idée.} Regrouper les 13 wilayas en $k$ classes homogènes selon leurs
8 indicateurs normalisés. Deux wilayas d'une même classe ont des profils
d'exclusion proches $\Rightarrow$ la même stratégie d'investissement.

\begin{meth}{Algorithme et objectif}
On minimise l'\emph{inertie} (somme des distances quadratiques intra-classe) :
\[
J \;=\; \sum_{c=1}^{k}\sum_{i\in C_c}\lVert x_i - \mu_c \rVert^2 .
\]
Algorithme itératif (Lloyd) :
\begin{enumerate}[leftmargin=2em]
  \item initialiser $k$ centres $\mu_c$ aléatoirement ;
  \item \textbf{affectation} : $C_c \gets \{i : c=\arg\min_c \lVert x_i-\mu_c\rVert\}$ ;
  \item \textbf{mise à jour} : $\mu_c \gets \frac{1}{|C_c|}\sum_{i\in C_c} x_i$ ;
  \item répéter 2--3 jusqu'à convergence (inertie stable).
\end{enumerate}
\end{meth}

\begin{where}
\txt{backend/analytics/clustering.py} : \txt{KMeans(n\_clusters=k, n\_init=20)}
sur données z-scorées. $k=3$ retenu pour une typologie lisible (modérée /
dominante « mahadra » / dominante « aucune instruction »). Endpoint
\txt{/api/clusters}, pages \txt{/clusters}, \txt{/wilayas} et \txt{/rapport}.
\end{where}

% ----------------------------------------------------------------------
\subsection{Choisir $k$ : inertie (coude) et silhouette}
\label{sec:silhouette}

\textbf{Idée.} L'inertie $J(k)$ décroît toujours avec $k$ ; le bon $k$ est
celui où la courbe « casse » (méthode du coude). La \emph{silhouette} mesure
la cohésion d'un point avec sa classe par rapport à la classe voisine.

\begin{meth}{Silhouette}
Pour chaque point $i$ :
\begin{align*}
a(i) &= \text{distance moyenne aux points de sa propre classe},\\
b(i) &= \min\ \text{des distances moyennes aux autres classes},\\
s(i) &= \frac{b(i)-a(i)}{\max\big(a(i),\,b(i)\big)} \;\in\; [-1,\,1].
\end{align*}
$s(i)\approx 1$ : point bien classé ; $s(i)<0$ : point probablement
mal affecté. On choisit le $k$ qui maximise la \emph{silhouette moyenne}.
\end{meth}

\begin{where}
\txt{backend/analytics/clustering.py} -- \txt{choose\_k()} balaye $k\in[2,6]$,
calcule inerties + \txt{silhouette\_score} ; le $k$ optimal est ensuite
raffiné pour la lisibilité (voir la note dans \txt{run\_all.py}).
\end{where}

% ======================================================================
\section{Apprentissage machine supervisé : les déterminants individuels}
% ======================================================================

% ----------------------------------------------------------------------
\subsection{Régression logistique}
\label{sec:logit}

\textbf{Idée.} Prédire une variable \emph{binaire} -- \txt{hors\_ecole}
(0/1) pour chaque enfant de 6--14 ans de l'EPCV (16\,451 enfants) -- en
fonction de facteurs : sexe, âge, milieu, pauvreté, éducation du chef de
ménage, wilaya.

\begin{meth}{Modèle}
La probabilité d'être hors école s'écrit avec la fonction logistique :
\[
p \;=\; P(y=1\mid x) \;=\; \frac{1}{1+e^{-(\beta_0 + \beta_1 x_1 + \cdots + \beta_p x_p)}}.
\]
Ou, en logit (log des cotes) :
\[
\log\frac{p}{1-p} \;=\; \beta_0 + \beta_1 x_1 + \cdots + \beta_p x_p.
\]
Les $\beta$ sont estimés par \textbf{maximum de vraisemblance} : on maximise
$L(\beta)=\prod_i p_i^{y_i}(1-p_i)^{1-y_i}$.
\end{meth}

\begin{meth}{Odds-ratio et intervalle de confiance}
Le \emph{odds-ratio} de la variable $x_j$ est
\[
OR_j = e^{\beta_j},
\]
avec intervalle de confiance à 95\,\% $\big[e^{\beta_j - 1{,}96\,\mathrm{SE}_j},\,
e^{\beta_j + 1{,}96\,\mathrm{SE}_j}\big]$.
\begin{itemize}[leftmargin=2em]
  \item $OR>1$ : le facteur \emph{augmente} la probabilité d'exclusion ;
  \item $OR<1$ : il la \emph{réduit} ;
  \item tous les effets sont \emph{ajustés} : « toutes choses égales par ailleurs ».
\end{itemize}
\end{meth}

\begin{where}
\txt{backend/analytics/logit.py} : \txt{sm.Logit(...)} de
\texttt{statsmodels}, avec variables binaires (références : garçon,
10--14 ans, urbain, non pauvre, Nouakchott). Résultats : \txt{odds\_ratio},
IC, \txt{pvalue}, \txt{pseudo\_r2}, \txt{aic}. Endpoint \txt{/api/logit},
affiché sur \txt{/} (Home) et \txt{/rapport}. Exemple de lecture : un enfant en milieu rural a
un $OR>1$ par rapport à l'urbain, \emph{à sexe, âge, pauvreté et wilaya
égaux}.
\end{where}

\begin{meth}{Variable écartée et pseudo-R²}
La variable \txt{nivcm}, étiquetée « Niveau d'éducation du CM » dans le
fichier SPSS, n'est \textbf{pas} l'éducation parentale : le croisement
\txt{C4N}$\,\times\,$\txt{nivcm} sur les 60\,600 individus est exactement
bloc-diagonal, c'est-à-dire un simple regroupement du niveau d'instruction
\emph{de l'individu lui-même}. L'introduire dans le modèle reviendrait à
régresser « hors école » sur « n'a aucun niveau scolaire » : une tautologie,
qui explique la séparation complète observée. Elle est donc écartée du logit
comme des règles d'association. Le fichier fourni ne contient aucune variable
d'éducation parentale, et cet effet n'est pas estimable ici.
\[
\text{pseudo-}R^2_{\text{McFadden}} = 1 - \frac{\ln L(\text{modèle complet})}
{\ln L(\text{modèle constant})}, \qquad
\text{AIC} = 2p - 2\ln L.
\]
\end{meth}

% ======================================================================
\section{Data mining : règles d'association (Apriori)}
\label{sec:apriori}

\textbf{Idée.} Trouver des \emph{combinaisons de facteurs} fréquemment
associées au hors-école, sans hypothèse a priori : âge, sexe, milieu rural,
pauvreté, région $\to$ hors école.

\begin{meth}{Support, confiance, lift}
Pour une règle $A \Rightarrow B$ (ici $B=\{\text{hors école}\}$) :
\[
\text{support}(A\Rightarrow B) = \frac{\text{freq}(A\cup B)}{N},
\qquad
\text{confiance}(A\Rightarrow B) = \frac{\text{support}(A\cup B)}{\text{support}(A)},
\]
\[
\text{lift}(A\Rightarrow B) = \frac{\text{confiance}(A\Rightarrow B)}{\text{support}(B)}.
\]
\begin{itemize}[leftmargin=2em]
  \item \emph{Support} : fréquence de la combinaison dans la population ;
  \item \emph{Confiance} : probabilité de $B$ sachant $A$ ;
  \item \emph{Lift} : rapport à l'indépendance. $\text{lift}>1$ : $A$ rend
        $B$ \emph{plus probable} qu'au hasard ; c'est le critère de
        pertinence utilisé ici (seuil $>1{,}05$).
\end{itemize}
\end{meth}

\begin{meth}{Algorithme Apriori}
\begin{enumerate}[leftmargin=2em]
  \item générer les \emph{itemsets} de taille $1$ dont le support dépasse
        $\min\_support$ ;
  \item construire itérativement les itemsets de taille $k+1$ par
        \emph{jointure} d'itemsets de taille $k$ (propriété d'anti-monotonie :
        tout sur-ensemble d'un itemset non fréquent est non fréquent) ;
  \item produire les règles et ne garder que celles satisfaisant
        $\min\_confidence$ et un lift $>1$.
\end{enumerate}
\end{meth}

\begin{where}
\txt{backend/analytics/rules.py} : \txt{apriori()} et
\txt{association\_rules()} de \texttt{mlxtend} sur les 16\,451 enfants.
Filtres : conséquence = exactement \{\txt{hors\_ecole}\}, lift $>1{,}05$.
Endpoints : \txt{/api/rules}.
\textbf{Page utilisée} : \txt{/regles} (page dédiée du frontend), qui
affiche les règles triées par lift, les profils socio-démographiques les plus
sûrs (\emph{top\_confiance}) et les règles par région (\emph{par\_region}).
Résultat : 21 règles non triviales, la plus forte étant
\{rural, 6--9 ans, Sud-Est\} $\Rightarrow$ hors école, à
\textbf{69{,}5\,\% de confiance et lift 2{,}10}. Aucune règle n'atteint
100\,\% : les règles déterministes de la version précédente provenaient
uniquement de la variable tautologique \txt{nivcm} (cf.\ §~\ref{sec:logit}).
\end{where}

% ======================================================================
\section{Analyse de graphes et détection de communautés}
\label{sec:graphs}

% ----------------------------------------------------------------------
\subsection{Représentation en graphe}

\textbf{Idée.} Un \emph{graphe} $G=(V,E)$ est un ensemble de nœuds $V$
(wilayas ou indicateurs) reliés par des arêtes $E$ pondérées (par $|r|$).
La structure du graphe révèle des groupes, des hubs et des isolés invisibles
dans un tableau.

\begin{meth}{Degré et matrice d'adjacence}
Le \emph{degré} $\deg(v)$ d'un nœud est le nombre de ses voisins.
La \emph{matrice d'adjacence} $A$ vaut $A_{uv}=1$ si $(u,v)\in E$, et la
matrice \emph{pondérée} $W_{uv}=|r_{uv}|$ sinon $0$. Le degré (pondéré) est
$\deg(v)=\sum_u W_{uv}$. Dans EduFocus, le degré des indicateurs détermine
la \emph{taille des nœuds} du graphe.
\end{meth}

\begin{where}
\txt{backend/analytics/graph.py} :
\begin{itemize}[leftmargin=2em]
  \item \textbf{Réseau de similarité des wilayas} : nœuds = wilayas, arêtes =
        paires dont $|r|\ge0{,}85$ (top-2 voisins), poids = $|r|$.
  \item \textbf{Graphe de corrélations des indicateurs} : nœuds =
        indicateurs, arêtes si $|r|\ge0{,}75$, taille des nœuds $\propto$
        degré. Endpoints \txt{/api/graph/similarite} et
        \txt{/api/graph/correlations} ; page \txt{/reseau}.
\end{itemize}
\end{where}

% ----------------------------------------------------------------------
\subsection{Détection de communautés : Louvain}

\textbf{Idée.} Regrouper les wilayas qui partagent le même profil
d'exclusion. La méthode de \emph{Louvain} maximise la \emph{modularité} $Q$.

\begin{meth}{Modularité de Newman--Girvan}
\[
Q \;=\; \frac{1}{2m}\sum_{u,v}
\Big[\,W_{uv} - \frac{d_u d_v}{2m}\Big]\,\delta(c_u,c_v),
\]
avec $m=\frac12\sum_{u,v}W_{uv}$ (poids total des arêtes), $d_u$ le degré
pondéré de $u$, et $\delta(c_u,c_v)=1$ si $u$ et $v$ sont dans la même
communauté. $Q\in[-1,1]$ ; $Q>0{,}3$ indique une structure communautaire
réelle (meilleure que le hasard).
\end{meth}

\begin{meth}{Algorithme de Louvain}
\begin{enumerate}[leftmargin=2em]
  \item \textbf{Phase locale} : chaque nœud est déplacé vers la communauté
        voisine qui augmente le plus $Q$ (gain de modularité), jusqu'à
        stabilisation ;
  \item \textbf{Agrégation} : les communautés deviennent des super-nœuds ;
  \item on répète 1--2 jusqu'à convergence.
\end{enumerate}
\end{meth}

\begin{where}
\txt{backend/analytics/graph.py} :
\txt{community\_louvain.best\_partition(G, random\_state=42)} (python-louvain).
Résultat : \textbf{3 communautés} de wilayas au profil proche, exposées sur
\txt{/reseau} (couleurs des nœuds) et dans le rapport PDF \txt{/rapport}.
\end{where}

% ----------------------------------------------------------------------
\subsection{Layout « force-directed » (carte du graphe)}

\textbf{Idée.} Positionner les nœuds dans le plan en simulant un système de
forces physique : les arêtes tirent les nœuds (ressorts), les nœuds se
repoussent (répulsion de Coulomb).

\begin{meth}{Force dirigée (modèle ressort)}
Force de ressort entre nœuds reliés : $F_s = k_s\,(d - l_0)$ (Hooke) ;
répulsion entre tous les nœuds : $F_r = \dfrac{k_r}{d^2}$ ; amortissement
$\eta$ : on itère $v \mathrel{+}= (F_s+F_r)/m$, $x \mathrel{+}= \eta\,v$.
La position finale est une lecture visuelle de la structure : wilayas
connectées se rapprochent, communautés forment des groupes visibles.
\end{meth}

\begin{where}
Frontend : \txt{frontend/src/lib/charts.ts} (\txt{forceGraph()}) avec
ECharts (layout \emph{force}). Taille des nœuds wilayas $\propto$ nombre
d'enfants hors école ; page \txt{/reseau}.
\end{where}

% ======================================================================
\section{Indicateur composite : l'Indice de Priorité Éducative (IPE)}
\label{sec:ipe}

\textbf{Idée.} Combiner trois dimensions en \emph{une} note 0--100 pour
classer les wilayas : volume du problème, intensité de l'exclusion,
vulnérabilité du territoire.

\begin{meth}{Formule}
\[
\boxed{\;\text{IPE} \;=\; 0{,}40\,\text{Volume} + 0{,}35\,\text{Intensité}
+ 0{,}25\,\text{Vulnérabilité}\;}
\]
avec (toutes normalisées en 0--100, voir §~\ref{sec:minmax}) :
\begin{itemize}[leftmargin=2em]
  \item $\text{Volume} = \log(1+\text{enfants hors école})$ normalisé ;
  \item $\text{Intensité} = \text{taux hors école (\%)}$ normalisé ;
  \item $\text{Vulnérabilité} = 0{,}75\,\text{pauvreté} + 0{,}25\,\text{dépendance jeunes}$.
\end{itemize}
Le rang est la \emph{position décroissante} de l'IPE sur les 13 wilayas
(méthode \txt{min} pour les ex æquo).
\end{meth}

\begin{where}
\txt{backend/analytics/index.py} (\txt{build\_index()}) : classe les wilayas
et dérive un \emph{levier d'action} dominant par règles de seuil (construire
des écoles / passerelles mahadra / rétention). Endpoints \txt{/api/summary}
et \txt{/api/wilayas} ; utilisé sur \txt{/}, \txt{/wilayas}, \txt{/carte}
(choroplèthe) et \txt{/rapport}.
\end{where}

% ======================================================================
\section{Scénarios prospectifs à 2030}
\label{sec:scenarios}

\textbf{Idée.} Traduire les leviers en trajectoires chiffrées : tendance
« laisser-faire », passerelles mahadra $\to$ formel (25\,\% et 50\,\%),
construction d'écoles de proximité.

\begin{meth}{Projection}
Avec $\tau_0=33{,}1\,\%$ (2022), pente annuelle $\beta$ (§~\ref{sec:ols}) et
$n_y = 2030-2022$ :
\[
\tau_{2030} = \max\big(\tau_0 + \beta\,n_y,\;0\big)
\]
puis, si l'on convertit une part $s$ des $M_0$ enfants mahadra (199\,493)
vers le formel :
\[
\tau_{2030} \leftarrow \tau_{2030} - \frac{s\,M_0}{P_{6-14}} \times 100
\]
où $P_{6-14}=1\,135\,657$. Besoin en écoles : on cible une densité de 1
établissement pour 1\,000 enfants,
$\text{besoin}_w = \max\big(1 - d_w,\,0\big)\times \frac{\text{pop}_{6-14}^{(w)}}{1000}$,
et coût estimé à 25\,M\,MRO par établissement.
\end{meth}

\begin{where}
\txt{backend/analytics/scenarios.py} : 4 scénarios comparés (taux et
effectifs 2030, réduction vs 2022, coûts). Endpoint \txt{/api/scenarios},
page \txt{/strategies} et rapport \txt{/rapport}.
\end{where}

% ======================================================================
\section{Récapitulatif : méthode $\leftrightarrow$ où c'est utilisé}
\label{sec:resume}
% ======================================================================

\begin{table}[h!]
\centering
\footnotesize
\begin{tabularx}{\textwidth}{@{}l l l X X@{}}
\toprule
\textbf{Famille} & \textbf{Méthode} & \textbf{Backend} & \textbf{Endpoint} & \textbf{Pages} \\
\midrule
Stat. & Normalisation min--max & \txt{index.py} & \txt{/api/wilayas} & \txt{/}, \txt{/wilayas}, \txt{/carte}, \txt{/rapport} \\
Stat. & z-score & \txt{clustering.py} & \txt{/api/clusters} & \txt{/clusters}, \txt{/wilayas}, \txt{/rapport} \\
Stat. & log1p & \txt{index.py}, \txt{matrice.py} & \txt{/api/wilayas}, \txt{/api/matrice} & \txt{/strategies} \\
Stat. & Médianes (quadrants) & \txt{matrice.py} & \txt{/api/matrice} & \txt{/strategies}, \txt{/rapport} \\
Stat. & Corrélation de Pearson & \txt{indicateurs.py}, \txt{graph.py} & \txt{/api/indicateurs}, \txt{/api/graph/*} & \txt{/indicateurs}, \txt{/reseau}, \txt{/rapport} \\
Stat. & Pareto / Lorenz / Gini & \txt{concentration.py} & \txt{/api/concentration} & \txt{/}, \txt{/strategies}, \txt{/rapport} \\
Stat. & OLS (tendance) & \txt{scenarios.py} & \txt{/api/scenarios} & \txt{/strategies}, \txt{/rapport} \\
ML non sup. & K-means + silhouette & \txt{clustering.py} & \txt{/api/clusters} & \txt{/clusters}, \txt{/wilayas}, \txt{/rapport} \\
ML sup. & Régression logistique & \txt{logit.py} & \txt{/api/logit} & \txt{/}, \txt{/rapport} \\
DM & Apriori (règles) & \txt{rules.py} & \txt{/api/rules} & \txt{/regles} \\
Graphe & Louvain (communautés) & \txt{graph.py} & \txt{/api/graph/similarite} & \txt{/reseau} \\
Graphe & Force-directed (ECharts) & --- & --- & \txt{/reseau} \\
Index & IPE composite & \txt{index.py} & \txt{/api/summary} & \txt{/}, \txt{/wilayas}, \txt{/carte}, \txt{/rapport} \\
Prospectif & Projections 2030 & \txt{scenarios.py} & \txt{/api/scenarios} & \txt{/strategies}, \txt{/rapport} \\
\bottomrule
\end{tabularx}
\caption{Correspondance complète méthode $\to$ code $\to$ API $\to$ interface.}
\end{table}

\begin{notebox}
\textbf{Lecture croisée.} Les méthodes ne sont pas isolées : la normalisation
alimente le clustering et l'IPE ; la corrélation alimente les graphes ; les
règles d'association compensent la limite de la régression (séparation
complète). L'\textbf{IPE} hiérarchise, le \textbf{clustering} regroupe,
les \textbf{graphes} révèlent la structure, la \textbf{logistique} explique
les causes individuelles, et les \textbf{scénarios} chiffrent les actions.
\end{notebox}

% ======================================================================
\section{Bibliothèque logicielle}
% ======================================================================

\begin{table}[h!]
\centering
\small
\begin{tabularx}{\textwidth}{@{}l l X@{}}
\toprule
\textbf{Bibliothèque} & \textbf{Rôle} & \textbf{Méthodes mobilisées} \\
\midrule
\texttt{numpy} / \texttt{pandas} & calcul & vecteurs, agrégats, tableaux \\
\texttt{scipy} & statistiques & linkage (CAH), distributions \\
\texttt{scikit-learn} & ML & \texttt{KMeans}, \texttt{StandardScaler},
\texttt{silhouette\_score} \\
\texttt{statsmodels} & stats & \texttt{Logit} (régression logistique, IC, p) \\
\texttt{networkx} & graphes & \texttt{Graph}, \texttt{degree} \\
\texttt{python-louvain} & graphes & \texttt{best\_partition} (Louvain) \\
\texttt{mlxtend} & data mining & \texttt{apriori}, \texttt{association\_rules} \\
\texttt{pyreadstat} & données & lecture EPCV 2019 (\texttt{.sav}) \\
\texttt{fastapi} / \texttt{uvicorn} & API & exposition des résultats \\
\bottomrule
\end{tabularx}
\caption{Bibliothèques Python utilisées par le backend (\txt{backend/requirements.txt}).}
\end{table}

\vspace{0.8cm}
\begin{center}
\small\color{eduMute}
EduFocus -- Datathon IndabaX Mauritanie 2026. Données : EPCV 2019 (ONS),
World Bank (WDI), 13 wilayas.
\end{center}

\end{document}
